\section{ISR 与缓存：把后端工作 materialize 到边缘}

电商 Black Friday 的典型瓶颈不是边缘网络，而是 origin/backend 无法承受每个用户都触发动态渲染。Incremental Static Regeneration（ISR）把一次 render 的结果 \term{materialize}（物化）为可缓存 artifact，在 freshness policy 到期或事件触发时重新生成。它改变了负载模型：后端处理“内容变化”，edge 处理“访问次数”。

\subsection{从 request amplification 到 backend shielding}

若页面在一个 freshness window 中收到 $R$ 次请求，传统动态渲染需要约 $R$ 次 origin work；理想 ISR 只需一次生成和一次 revalidation，后端工作量从 $O(R)$ 接近 $O(1)$。真实系统还要处理 stampede、失败回退、区域 cache、tag invalidation 和 stale content。

\lecturefigure{05-isr-cache-loop.png}{ISR 用 materialization 和 revalidation 把重复读取从 backend 移到 edge。}{本地字幕 05:55--07:20；Next.js 官方 ISR 文档。}

\begin{knowledgebox}{读图：比较访问频率与变化频率}
ISR 最适合“读很多、写较少、允许明确 freshness”的内容。若数据每次请求都不同，cache 命中率低；若内容绝不能陈旧，就需要同步 origin 或更严格 invalidation。核心问题不是“能否缓存”，而是业务允许什么 consistency contract。
\end{knowledgebox}

\begin{warningbox}{缓存不会让 origin 无限扩展}
Cache miss、同时 revalidation、个性化请求和大量不同 key 仍会击穿后端。平台需要 request coalescing、stale-while-revalidate、负载保护和可见的 cache status，开发者也必须理解 cache key 与 invalidation 语义。
\end{warningbox}

\teachervoice{Rauch 的表达很有工程感：云“elastic”不代表任何组件都完美扩展。平台价值之一，是自动组合多级缓存和 materialization，把开发者过去手工管理的 Memcached、worker pool 和横向扩展模式固化为框架能力。}

\subsection{本章小结}

ISR 把访问规模与后端计算解耦，让应用流量更多由 edge artifact 承担。它不是静态化一切，而是显式管理 freshness、materialization 和 revalidation。

